\documentclass[]{meetingminutes}
% Options:
%   language=catalan | spanish | english
%   novotechecks  (disables vote count vs. attendee count validation)
%   agenda        (generates an interactive agenda checklist document)
%   nosort        (disables alphabetical sorting globally)
%   nohierarchy   (flattens the attendee table without prioritizing roles)

% Font Package
\RequirePackage{tgheros}
\renewcommand{\familydefault}{\sfdefault}

% ==== Institution Metadata ===== 
\institution{LinuxUPC}
\department{Secretaría}
\headerlogo{./linux_logo.jpg}

% ==== Meeting Metadata ===== 
\title{TITOL}
\type{TIPUS_DE_REUNIO}
% Possible Types are:
% - GeneralAssembly  - BoardMeeting  - ExtraordinaryGeneralAssembly

\date{DATA}
\location{Despatx (A3S105)}
\starttime{START}
\endtime{END}

% ==== Attendees & Guests ===== 
% \attendee[chair]{Marcel Mula Tixier}
% \attendee[secretary]{Cristina Frias}
% \attendee[treasurer]{Chentao Chen}
% \attendee{Joan Jose}

% \guest[Afiliación]{Pepito Palotes}

\begin{document}
\maketitle

\agendaitem{Item 1}
\agendaitem{Item 2}
\agendaitem{Item 3}

\printtasks 
\end{document}

% Chuleta Bloques Especiales (v3.0)
% =================================
% IMPORTANTE: Las listas de personas en votaciones usan punto y coma (;) como separador.
%
% \vote{mocion}{lista a favor; ...}{lista en contra; ...}{lista abstencion; ...}  (¡Los guests no votan!)
% \anonymousvote{mocion}{cuantos a favor}{cuantos en contra}{cuantos abstencion}
% \election{mocion}{
%        \electionoption{opcion}{lista de votantes; ...}
%        \abstention{lista de abstenciones; ...}
% }
% \task[fecha_opcional]{responsable}{tarea}
